\Askhsh{38923}{4}
\begin{ylh}{2}
    \item 2.2. Διάταξη Πραγματικών Αριθμών
    \item 3.3. Εξισώσεις 2ου Βαθμού
\end{ylh}
Η Μαρία επιθυμεί να φτιάξει έναν ορθογώνιο κήπο έξω από το σπίτι της και, για να υπολογίσει το μήκος και το πλάτος του, θέλει να ικανοποιούνται δύο βασικές προδιαγραφές:

η περίμετρος του κήπου να είναι $26\ m$ και το εμβαδόν του να είναι $40{\ m}^{2}$.

\begin{alist}
\item Να κατασκευάσετε μια εξίσωση, η οποία να έχει ως λύσεις το μήκος και το πλάτος του κήπου.
\monades{6}
\item Ποιες είναι οι διαστάσεις του κήπου (μήκος και πλάτος);\monades{5}
\item Η Μαρία θέλει να βάλει στη μέση του κήπου ένα στρογγυλό παρτέρι με εμβαδόν $20m^{2}$. Θα χωρέσει το παρτέρι στον κήπο; Να αιτιολογήσετε την απάντησή σας.

(Δίνεται το εμβαδόν κύκλου $E = \pi\rho^{2},$ όπου $\rho$ η ακτίνα του κύκλου και $\pi \cong 3,14)$.
\monades{7}
\item Η Μαρία σκέφτεται μήπως αλλάξει το πλάτος του κήπου, διατηρώντας σταθερό το εμβαδόν του στα $40\ m^{2}$. Επειδή όμως περιορίζεται από την διαθέσιμη επιφάνεια του οικοπέδου του σπιτιού, η νέα διάσταση του πλάτους πρέπει να είναι μικρότερη από $7\ m$ και μεγαλύτερη από $4\ m\ .$ Να υπολογίσετε το διάστημα των τιμών που μπορεί να έχει το μήκος του κήπου.
\monades{7}

Δίνεται $\sqrt{6,37} \cong 2,52$.

{\footnotesize Το παραπάνω θέμα αναπτύχθηκε στο πλαίσιο του έργου: «Ανάπτυξη Δοκιμασιών Αξιολόγησης Δεξιοτήτων Εγγραμματισμού στα μαθήματα της Νεοελληνικής Γλώσσας και Λογοτεχνίας, της Άλγεβρας, της Φυσικής και της Χημείας Α' Λυκείου Γενικού Λυκείου» Ανάδοχος: «Ειδικός Λογαριασμός Κονδυλίων Έρευνας (Ε.Λ.Κ.Ε) Πανεπιστημίου Ιωαννίνων» ΑΔΑΜ: 25SYMV016348911 2025-02-20.}
\end{alist}

